% Options for packages loaded elsewhere
\PassOptionsToPackage{unicode}{hyperref}
\PassOptionsToPackage{hyphens}{url}
\PassOptionsToPackage{dvipsnames,svgnames,x11names}{xcolor}
\documentclass[
  10pt,
  fontset=fandol]{ctexart}
\usepackage{xcolor}
\usepackage[margin=16mm]{geometry}
\usepackage{amsmath,amssymb}
\setcounter{secnumdepth}{-\maxdimen} % remove section numbering
\usepackage{iftex}
\ifPDFTeX
  \usepackage[T1]{fontenc}
  \usepackage[utf8]{inputenc}
  \usepackage{textcomp} % provide euro and other symbols
\else % if luatex or xetex
  \usepackage{unicode-math} % this also loads fontspec
  \defaultfontfeatures{Scale=MatchLowercase}
  \defaultfontfeatures[\rmfamily]{Ligatures=TeX,Scale=1}
\fi
\usepackage{lmodern}
\ifPDFTeX\else
  % xetex/luatex font selection
\fi
% Use upquote if available, for straight quotes in verbatim environments
\IfFileExists{upquote.sty}{\usepackage{upquote}}{}
\IfFileExists{microtype.sty}{% use microtype if available
  \usepackage[]{microtype}
  \UseMicrotypeSet[protrusion]{basicmath} % disable protrusion for tt fonts
}{}
\makeatletter
\@ifundefined{KOMAClassName}{% if non-KOMA class
  \IfFileExists{parskip.sty}{%
    \usepackage{parskip}
  }{% else
    \setlength{\parindent}{0pt}
    \setlength{\parskip}{6pt plus 2pt minus 1pt}}
}{% if KOMA class
  \KOMAoptions{parskip=half}}
\makeatother
\setlength{\emergencystretch}{3em} % prevent overfull lines
\providecommand{\tightlist}{%
  \setlength{\itemsep}{0pt}\setlength{\parskip}{0pt}}
\usepackage{amsmath,amssymb,mathtools,needspace}
\xeCJKDeclareCharClass{CJK}{"2032,"0400->"04FF,"2160->"216B,"2460->"2473,"25A0,"25C7,"25CB,"25CF}
\IfFontExistsTF{Arial Unicode MS}{\xeCJKsetup{AutoFallBack=true}\setCJKfallbackfamilyfont{\CJKrmdefault}{Arial Unicode MS}}{}
\setcounter{secnumdepth}{0}
\setlength{\emergencystretch}{2em}
\allowdisplaybreaks
\usepackage{bookmark}
\IfFileExists{xurl.sty}{\usepackage{xurl}}{} % add URL line breaks if available
\urlstyle{same}
\hypersetup{
  pdftitle={习题},
  colorlinks=true,
  linkcolor={Maroon},
  filecolor={Maroon},
  citecolor={Blue},
  urlcolor={Blue},
  pdfcreator={LaTeX via pandoc}}

\title{习题}
\author{}
\date{}

\begin{document}
\maketitle

\subsubsection{习题}\label{ux4e60ux9898}

\begin{enumerate}
\def\labelenumi{\arabic{enumi}.}
\item
  设 \(x>0, x\) 的相对误差为 \(\delta\), 求 \(\ln x\) 的误差.
\item
  设 \(x\) 的相对误差为 \(2\%\), 求 \(x^n\) 的相对误差.
\item
  下列各数都是经过四舍五入得到的近似数, 即误差限不超过最后一位的半个单位, 试指出它们是几位有效数字:
\end{enumerate}

\[
x_1^* = 1.1021, x_2^* = 0.031, x_3^* = 385.6, x_4^* = 56.430, x_5^* = 7 \times 1.0.
\]

\begin{enumerate}
\def\labelenumi{\arabic{enumi}.}
\setcounter{enumi}{3}
\tightlist
\item
  利用式(1.3.3)求下列近似值的误差限:
\end{enumerate}

\begin{enumerate}
\def\labelenumi{(\arabic{enumi})}
\tightlist
\item
  \(x_1^* + x_2^* + x_4^*\), (2) \(x_1^* x_2^* x_3^*\), (3) \(x_2^* / x_4^*\), 其中 \(x_1^*, x_2^*, x_3^*, x_4^*\) 均为习题 3 所给的数.
\end{enumerate}

\begin{enumerate}
\def\labelenumi{\arabic{enumi}.}
\setcounter{enumi}{4}
\item
  计算球体积要使相对误差限为 \(1\%\), 问度量半径为 \(R\) 时允许的相对误差限是多少.
\item
  设 \(Y_0=28\), 按递推公式
\end{enumerate}

\[
Y_n = Y_{n-1} - \frac{1}{100} \sqrt{783} \quad (n=1,2,\cdots)
\]

计算到 \(Y_{100}\). 若取 \(\sqrt{783} \approx 27.982\) (五位有效数字), 试问计算 \(Y_{100}\) 将有多大误差.

\begin{enumerate}
\def\labelenumi{\arabic{enumi}.}
\setcounter{enumi}{6}
\item
  求方程 \(x^2-56x+1=0\) 的两个根, 使它至少具有四位有效数字 (\(\sqrt{783} \approx 27.982\)).
\item
  当 \(N\) 充分大时, 怎样求 \(\int_N^{N+1} \frac{1}{1+x^2} \mathrm{d}x\)?
\item
  正方形的边长大约为 \(100 \mathrm{~cm}\), 应怎样测量才能使其面积误差不超过 \(1 \mathrm{~cm}^2\)?
\item
  设 \(s=\frac{1}{2} g t^2\), 假定 \(g\) 是准确的, 而对 \(t\) 的测量有 \(\pm 0.1 \mathrm{~s}\) 的误差, 证明: 当 \(t\) 增加时 \(s\) 的绝对误差增大, 而相对误差却减小.
\item
  序列 \(\{y_n\}\) 满足递推关系 \(y_n=10y_{n-1}-1 (n=1,2,\cdots)\), 若 \(y_0=\sqrt{2} \approx 1.41\) (三位有效数字),
\end{enumerate}

\begin{quote}
校注（agent 补充）：习题3第五个数的原图清楚写为 \(x_5^*=7\times1.0\)，未作简化或改写。习题11在本页仅给出条件，题目后半句续下页。
\end{quote}

\href{../../source-images/pdf-026.jpeg}{查看原页}

计算到 \(y_{10}\) 时误差有多大？这个计算过程稳定吗？

\begin{enumerate}
\def\labelenumi{\arabic{enumi}.}
\setcounter{enumi}{11}
\tightlist
\item
  计算 \((\sqrt{2}-1)^6\), 取 \(\sqrt{2} \approx 1.4\), 利用下式计算, 哪一个得到的结果最好?
\end{enumerate}

\[
\frac{1}{(\sqrt{2}+1)^6}, \quad (3-2\sqrt{2})^3, \quad \frac{1}{(3+2\sqrt{2})^3}, \quad 99-70\sqrt{2}.
\]

\begin{enumerate}
\def\labelenumi{\arabic{enumi}.}
\setcounter{enumi}{12}
\tightlist
\item
  \(f(x)=\ln(x-\sqrt{x^2-1})\), 求 \(f(30)\) 的值. 求对数时误差有多大? 若改用另一等价公式
\end{enumerate}

\[
\ln(x-\sqrt{x^2-1})=-\ln(x+\sqrt{x^2+1})
\]

计算, 求对数时误差有多大?

\begin{enumerate}
\def\labelenumi{\arabic{enumi}.}
\setcounter{enumi}{13}
\item
  试用消元法解方程组 \(\begin{cases} x_1+10^{10}x_2=10^{10}, \\ x_1+x_2=2, \end{cases}\) 假定只用三位数计算, 结果是否可靠?
\item
  已知三角形面积 \(S=\frac{1}{2}ab\sin c\), 其中 \(c\) 为弧度, \(0<c<\frac{\pi}{2}\), 且测量 \(a, b, c\) 的误差分别为 \(\Delta a\), \(\Delta b\), \(\Delta c\). 证明面积的误差 \(\Delta S\) 满足
\end{enumerate}

\[
\left|\frac{\Delta S}{S}\right| \leqslant \left|\frac{\Delta a}{a}\right| + \left|\frac{\Delta b}{b}\right| + \left|\frac{\Delta c}{c}\right|.
\]

\begin{quote}
校注（agent 补充，CH01-E007）：习题13所谓``等价公式''的右端，原图清楚印为 \(-\ln(x+\sqrt{x^2+1})\)，此处保留加号。对于题中实值范围，\((x-\sqrt{x^2-1})(x+\sqrt{x^2-1})=1\)，故相应等价式的右端根号内应为 \(x^2-1\)。这是原书符号疑误，不是 OCR 错误；本页只转录题目，未代做习题。
\end{quote}

\begin{quote}
校注（agent 补充，CH01-E008）：习题15把面积的有限误差直接写成上述不等式，原文照录。该式可按一阶误差或微分估计理解；若 \(\Delta\) 表示有限误差，乘积的高阶项不能一般地省略。例如 \(\Delta a/a=\Delta b/b=0.01\) 且 \(\Delta c=0\) 时，精确相对面积变化为 \(1.01^2-1=0.0201>0.02\)。因此题目作为严格有限误差上界缺少限定。
\end{quote}

\subsection{本节覆盖原页的校注记录}\label{ux672cux8282ux8986ux76d6ux539fux9875ux7684ux6821ux6ce8ux8bb0ux5f55}

下列记录按原页关联，含同页相邻小节的校注；计算时同时读取上文校注和完整原页。

\begin{itemize}
\tightlist
\item
  \texttt{CH01-E007}：原书 PDF 页 26
\item
  \texttt{CH01-E008}：原书 PDF 页 26
\end{itemize}

\end{document}
